\documentclass[11pt]{article}
\usepackage{amsmath,amssymb,hyperref}
\title{Presentations and Restricted Descriptions: An Exploratory Note}
\author{Luca Blanchi}
\date{November 6, 2026}
% This preliminary article was generated or revised with ChatGPT (GPT-5.5+).
% It has not been fully verified or peer-reviewed and may contain errors or gaps.
\begin{document}
\maketitle
\begin{abstract}
A description of a mathematical object is useful only relative to specified
rules for interpreting it. Two elementary examples show how those rules affect
the shortest admissible description. They motivate a modelling question, but do
not establish a new general theory or a new mathematical result. This note
replaces a larger collection whose claims did not survive critical review in
publishable form.
\end{abstract}

\section{The question and its limits}
A group presentation specifies generators and relations, with the presented
group obtained by a fixed construction. One can ask a broader question: does
restricting how an object may be described make a useful mathematical problem
visible? This question occurs in many established settings, including generating
sets, normal forms and data compression. A common vocabulary does not by itself
supply a common theorem.

At a minimum, fix a collection $P$ of descriptions, a collection $X$ of objects,
and an interpretation map $R:P\to X$. Two descriptions are equivalent when their
images agree. If objects are considered only up to isomorphism, use isomorphism
classes in $X$, or specify that equivalence separately. A cost $c:P\to[0,\infty]$
defines
\[
C(x)=\inf_{\substack{p\in P\\R(p)=x}}c(p).
\]
This is a definition, not a complexity theorem. The infimum of the empty set is
$\infty$; a finite infimum need not be attained. For instance, descriptions of a
single object with costs $1/n$, for positive integers $n$, have infimum zero and
no cheapest description. Counting generators, instructions or encoded bits gives
different questions. The interpretation rule and its computational cost must
also be stated.

\section{A small example with inference rules}
Let $D=\{0,a,b,c\}$, where
\[
\begin{aligned}
0&=(0,0),& a&=(1,0),\\
b&=(0,1),& c&=(1,1),
\end{aligned}
\]
and let $R$ be coordinatewise order on $D$. A description is a subset $T$ of the
strict pairs of $R$. Its interpretation is the relation generated by specified
inference rules. Its cost is $|T|$; the points of $D$ and the rules are fixed in
advance and are not charged to this cost.

First allow only reflexivity and transitivity. To generate $R$ one needs all four
pairs
\[
(0,a),\quad (0,b),\quad (a,c),\quad (b,c).
\]
Each is a cover relation: there is no intermediate point through which it could
be derived by transitivity. These four pairs suffice, since they also derive
$(0,c)$. Thus the minimum cost is four.

Now additionally allow the following rule. If $x$ and $y$ are related to the same
upper endpoint $z$, relate to $z$ every $w\in D$ whose coordinates satisfy
$w_i\in\{x_i,y_i\}$ for $i=1,2$. This rule preserves coordinatewise order: if
$x_i,y_i\le z_i$, then also $w_i\le z_i$.

The three pairs
\[
T=\{(0,a),(0,b),(0,c)\}
\]
now suffice. Mix the lower endpoints in $(0,c)$ and the reflexive pair $(c,c)$ to
obtain $(a,c)$ and $(b,c)$. All strict pairs of $R$ have then been generated.
Three is also a lower bound. For any endpoint $z$, a nonreflexive pair ending at
$z$ cannot be created unless some generator already ends at $z$: transitivity
would require such a pair as its final strict step, while mixing $(z,z)$ with
itself creates nothing. Each of $a,b,c$ therefore needs an incoming generator.

The minimum has changed from four to three because the interpretation rule has
changed. There is no encoding-independent compression claim here. A new rule has
been made available for free. The calculation is elementary and is included only
to make this dependence explicit; it is not claimed as a research contribution.
Nor does this autonomous rule imply a theorem about any larger geometric
constraint system.

\section{A small example with copying}
Consider finite words over a fixed alphabet. A literal instruction appends one
symbol. With only literals, a word of length $N$ costs $N$ instructions. Enlarge the language by allowing a
previous-copy instruction: at output position $t$, choose an earlier position
$s<t$ and a length $L$, then sequentially append the symbols at positions
$s,s+1,\ldots,s+L-1$. Positions are numbered starting at zero. Overlap is allowed:
each symbol is read after it has already been written, since the source remains
$t-s$ positions behind the output.

For $N\ge2$, the constant word $a^N$ has a two-instruction description: write $a$,
then copy from position zero for length $N-1$. One instruction cannot suffice,
because no previous source exists at the start. Thus its minimum is two in this
model. The decoder, alphabet and parameter conventions have been fixed; the
length parameter still requires storage. With an ordinary binary encoding this
description uses $O(\log N)$ bits, not a constant number of bits.

This familiar copying mechanism makes a restricted-description question
concrete. One may ask which words require many copy instructions even though a
short program can compute them. That is a question in compression and
computability, and it should be investigated using their existing methods. The
example alone provides no lower bound for a nonconstant word and no evidence of
an additional general theory.

\section{What this viewpoint contributes here}
In the first example, fixing the target relation and varying the inference rules
produces a precise minimum-generator question. In the second, fixing the word
and varying the instruction language separates phrase count from encoded length.
These comparisons are reproducible. They justify stating the representation
model before discussing its cost. Both arguments can be expressed completely in
standard terminology without a new formalism.

The reviewed material does not establish that a general notion of presentation
produced a new theorem. Some concrete arguments belong to established subjects;
others require corrections or unresolved external inputs. Their provenance does
not establish a discovery heuristic, and they are not evidence for one here.
In particular, a verifier for proposed certificates does not prove that suitable
certificates exist. A descent argument requires actual admissible steps and a
well-founded measure that decreases on them.

The conclusion is limited: retain these examples as an exploratory note, with
no claim of novelty or of a developed theory. A stronger claim would require a
specific result, a checked proof, comparison with the relevant literature, and a
concrete explanation of what the choice of descriptions contributes beyond a
change of terminology.
\end{document}
